\answer{OR16}
From every unabsorbed state, $a+b$ consecutive rightward steps force absorption. Such a block has probability $2^{-(a+b)}$. Hence $\tau$ has a geometric tail and all its moments are finite.

The symmetric independent increments make the following processes martingales:
\[
 S_n,\qquad S_n^2-n,\qquad S_n^3-3nS_n,
\]
\[
 S_n^4-6nS_n^2+3n^2+2n.
\]
These identities follow by expanding $(S_n\pm1)^j$ and averaging. Apply optional stopping first at the bounded time $\tau\wedge N$, then let $N\to\infty$. The stopped position is bounded in absolute value by $\max(a,b)$, and $\tau$ has finite moments, so dominated convergence justifies the limit.

Set $r=\Prob(S_\tau=b)$. The first two martingales give
\[
 br-a(1-r)=0,\qquad
 \E\tau=\E S_\tau^2=b^2r+a^2(1-r).
\]
Therefore
\[
 \boxed{r=\frac a{a+b}},\qquad \boxed{\E\tau=ab}.
\]
To obtain the second moment, write $A=\E[\tau\ind_{\{S_\tau=b\}}]$ and $B=\E[\tau\ind_{\{S_\tau=-a\}}]$. The cubic martingale yields
\[
 A+B=ab,\qquad
 bA-aB=\frac{\E S_\tau^3}{3}=\frac{ab(b-a)}3.
\]
Solving these two linear equations gives
\[
 A=\frac{ab(2a+b)}{3(a+b)},\qquad
 B=\frac{ab(a+2b)}{3(a+b)}.
\]
Consequently,
\[
 \E(\tau S_\tau^2)=b^2A+a^2B
       =\frac{ab(a^2+ab+b^2)}3,
\]
\[
 \E S_\tau^4=ab(a^2-ab+b^2).
\]
Substitute these values into the stopped quartic identity
\[
 0=\E S_\tau^4-6\E(\tau S_\tau^2)
       +3\E\tau^2+2\E\tau
\]
to obtain
\[
 \boxed{\E\tau^2=\frac{ab(a^2+3ab+b^2-2)}3}.
\]
